\subsection{示向约束与定位区域}

记 \(S_i=(x_i,y_i)^{\mathsf T}\)，单位方向向量为
\begin{equation}
  \vect u(\theta)=\bigl(\cos\theta,\sin\theta\bigr)^{\mathsf T},
\end{equation}
示向误差上界为 \(\varepsilon=1^{\circ}\)。定义二维叉积
\(\operatorname{cross}(\vect a,\vect b)=a_xb_y-a_yb_x\)，则第 \(i\) 次测向对应的角域为
\begin{equation}
  W_i=\left\{\vect x\in\R^2:
  \begin{aligned}
  &\operatorname{cross}\!\left(\vect u(\theta_i-\varepsilon),\vect x-S_i\right)\ge 0,\\
  &\operatorname{cross}\!\left(\vect u(\theta_i+\varepsilon),\vect x-S_i\right)\le 0
  \end{aligned}\right\}.
  \label{eq:wedge}
\end{equation}

目标圆域为 \(B=\{\vect x:\|\vect x\|_2\le1800\}\)，故定位区域定义为
\begin{equation}
  \Omega=B\cap\bigcap_{i=1}^{m}W_i.
  \label{eq:location-region}
\end{equation}
若求出的线性约束交集已包含在 \(B\) 内，则 \(\Omega\) 是凸多边形。若目标圆边界参与截断，应保留圆弧交点并在圆弧上额外检查极值；也可用足够密的正多边形近似目标圆，并报告离散误差上界。

\subsection{多边形直径与覆盖圆}

设凸多边形顶点按逆时针排序为 \(V=\{\vect v_1,\ldots,\vect v_n\}\)，则其直径为
\begin{equation}
  D(\Omega)=\max_{1\le i<j\le n}\|\vect v_i-\vect v_j\|_2.
  \label{eq:diameter}
\end{equation}
顶点数较少时可枚举顶点对；数据规模较大时可用旋转卡壳法将复杂度从 \(O(n^2)\) 降至 \(O(n)\)。

设一对直径端点为 \(A,B\)，中点 \(C=(A+B)/2\)，半径 \(r=D/2\)。直径圆覆盖定位区域的充要条件为
\begin{equation}
  \max_{\vect v\in V}\|\vect v-C\|_2\le r.
  \label{eq:coverage}
\end{equation}
因此，“直径等于区域直径”并不自动保证覆盖，必须按式~\eqref{eq:coverage} 对所有极点判定。
